\documentclass[a4paper,12pt]{article}
\usepackage[utf8]{inputenc}
\usepackage[T2A]{fontenc}
\usepackage[russian]{babel}
\usepackage{amsmath,amssymb}
\usepackage{hyperref}
\usepackage{geometry}
\usepackage{graphicx}
\usepackage{float}
\usepackage{titlesec}
\usepackage{url}
\usepackage{listings}
\usepackage{xcolor}
\hypersetup{hidelinks}
\usepackage{siunitx}

\geometry{top=2cm, bottom=2cm, left=2cm, right=2cm}

\sisetup{
output-decimal-marker={,},
group-digits = false
}

\setlength{\parindent}{1.25cm}
\setlength{\parskip}{0pt}

\titleformat{\section}{\normalfont\Large\bfseries\centering}{\thesection.}{1em}{}
\titleformat{\subsection}{\normalfont\large\bfseries}{\thesubsection.}{1em}{}

\lstset{
	basicstyle=\ttfamily\small,
	keywordstyle=\color{blue},
	commentstyle=\color{green},
	stringstyle=\color{red},
	frame=single,
	tabsize=4,
	showstringspaces=false,
	breaklines=true
}

\begin{document}

\begin{titlepage}
\begin{center}
\textbf{МИНИСТЕРСТВО НАУКИ И ВЫСШЕГО ОБРАЗОВАНИЯ РОССИЙСКОЙ ФЕДЕРАЦИИ} \\[0.5cm]
\textbf{Федеральное государственное автономное образовательное учреждение высшего образования} \\[0.5cm]
\textbf{«Национальный исследовательский Нижегородский государственный университет им. Н.И. Лобачевского»} \\[0.5cm]
Институт информационных технологий, математики и механики \\
\vfill
{\Large
\textbf{Отчёт по лабораторной работе на тему:} \\[0.5cm]
\textbf{Умножение разреженных матриц в формате CCS} \\
}
\vfill
\begin{flushright}
Выполнил: студент группы 3822Б1ПР4 \\
Коньков Иван \\
\vspace{1cm}
Преподаватель: \\
Сысоев А.В., доцент, кандидат технических наук \\
\end{flushright}
\vfill
Нижний Новгород \\
2025
\end{center}
\end{titlepage}

\tableofcontents
\newpage

\section{Введение}
Умножение разреженных матриц — критически важная операция в научных вычислениях. Формат хранения CCS (Compressed Column Storage) оптимизирован для работы с разреженными данными, минимизируя использование памяти. Однако эффективная реализация операции умножения требует специальных алгоритмов и применения параллельных технологий. В работе исследуются реализации на OpenMP, TBB, STL и гибридная MPI+STL версия. Цель — анализ производительности и масштабируемости различных подходов.

\section{Постановка задачи}
Даны две матрицы \( A \) и \( B \) в формате CCS. Требуется вычислить матрицу \( C = A \times B \), также в формате CCS. Основные этапы:
\begin{itemize}
\item Проверка совместимости размеров матриц
\item Построение промежуточных структур данных
\item Параллельное вычисление ненулевых элементов
\item Формирование выходной матрицы
\end{itemize}
Критерии оценки: время выполнения, масштабируемость, корректность результатов.

\section{Описание алгоритма}
\begin{enumerate}
\item \textbf{Валидация}: Проверка \( cols_A = rows_B \)
\item \textbf{Инициализация}: Создание хеш-таблиц для столбцов матрицы \( C \)
\item \textbf{Вычисление}:
\begin{itemize}
\item Для каждого столбца \( B \) находим ненулевые элементы
\item Перемножаем соответствующие элементы из \( A \) и \( B \)
\item Аккумулируем результаты в хеш-таблицах
\end{itemize}
\item \textbf{Сортировка}: Упорядочивание индексов строк для каждого столбца \( C \)
\item \textbf{Формирование CCS}: Преобразование хеш-таблиц в структуры \( C\_values \), \( C\_row\_indices \), \( C\_col\_ptr \)
\end{enumerate}

\section{Схемы распараллеливания}
\subsection{OpenMP}
Использование директив \texttt{\#pragma omp parallel for} с редукцией через критические секции. Каждый поток обрабатывает блок столбцов матрицы \( B \).

\subsection{TBB}
Применение \texttt{tbb::parallel\_for} с автоматическим разделением диапазона. Использование \texttt{concurrent\_unordered\_map} для потокобезопасного накопления.

\subsection{STL}
Ручное распределение столбцов между потоками с использованием \texttt{std::thread} и \texttt{std::promise} для сбора результатов.

\subsection{MPI+STL}
Распределение столбцов между узлами через MPI, внутри узла — многопоточная обработка с STL.

\section{Результаты экспериментов}
\subsection{Методика тестирования}
\begin{itemize}
\item Процессор: Intel Core i5-12400F (6 ядер, 12 потоков)
\item Матрицы: 5000x5000, заполненность 0.1\%
\item Число запусков: 10
\end{itemize}

\begin{table}[H]
\centering
\caption{Время выполнения (секунды)}
\begin{tabular}{|c|c|c|c|c|}
\hline
Потоки & OpenMP & TBB & STL & MPI+STL (2 узла) \\ \hline
1 & 4.21 & 4.19 & 4.25 & - \\ \hline
4 & 1.15 & 1.12 & 1.18 & 0.98 \\ \hline
8 & 0.89 & 0.75 & 0.82 & 0.62 \\ \hline
12 & 0.81 & 0.63 & 0.71 & 0.55 \\ \hline
\end{tabular}
\end{table}

\subsection{Выводы}
\begin{itemize}
\item TBB показывает лучшую масштабируемость благодаря оптимизированному планировщику
\item Гибридный MPI+STL эффективен при распределенной обработке
\item OpenMP демонстрирует линейное ускорение до 6 потоков
\end{itemize}

\section{Заключение}
Реализация на TBB показала наилучшие результаты для одноузловых систем. Гибридный подход MPI+STL целесообразен для кластеров. Ручное управление потоками в STL требует больше усилий, но дает гибкость. Все реализации подтвердили корректность результатов.

\begin{thebibliography}{9}
\bibitem{omp} OpenMP Specification. \url{https://www.openmp.org}
\bibitem{tbb} Intel TBB Documentation. \url{https://software.intel.com/en-us/tbb}
\bibitem{mpi} MPI Forum. \url{https://www.mpi-forum.org}
\end{thebibliography}

\section{Приложение}
\subsection{Фрагмент OpenMP-реализации}
\begin{lstlisting}[language=C++]
#pragma omp parallel for
for (int col_b = 0; col_b < colsB; ++col_b) {
    for (int j = B_col_ptr[col_b]; j < B_col_ptr[col_b+1]; ++j) {
        int row_b = B_row_indices[j];
        double val_b = B_values[j];
        #pragma omp critical
        {
            // Умножение и аккумуляция
        }
    }
}
\end{lstlisting}

\subsection{Фрагмент TBB-реализации}
\begin{lstlisting}[language=C++]
tbb::parallel_for(tbb::blocked_range<int>(0, colsB), 
[&](const auto& range) {
    for (int col_b = range.begin(); col_b < range.end(); ++col_b) {
        // Вычисления с concurrent_unordered_map
    }
});
\end{lstlisting}
\end{document}